\section{Changed terms and recovered credit}\label{sec:bank-restructuring}

A borrower can receive easier terms without recovering its ability to pay.
A bank can also recover its expectation of payment without recovering all
the credit quality present at origination. These are different changes to
the loan. The following historical Australian examples use the retained
AASB~9 compilation for synthetic 2024--2025 reporting periods. Probabilities,
legal findings and credit assessments are supplied assumptions; taxes and
fees are excluded.

Start from cash~100, a gross loan~110, allowance~55 and deposits~180. Net assets~155
leave equity $-25$. At the 2024 close, the bank replaces a payment~121 due
one year later with~110 at the same date. Assume the change preserves the
existing asset's contractual rights and does not qualify for derecognition.
The original annual effective interest rate is~10\%. Discounting the amended
contractual payment gives
\[
 G_{\mathrm{modified}}=\frac{110}{1.1}=100,
 \qquad \text{modification loss}=110-100=10.
\]
Debit modification expense~10 and credit gross loan~10. This calculation changes
the contractual carrying amount before considering expected defaults.
Paragraph~5.4.3 requires that separation. Fees would require an additional
carrying-value adjustment and amortisation method
\citep[paragraph~5.4.3]{aasb9dated}.

The bank separately expects full~110 with probability $4/5$ and zero with
probability $1/5$. Expected cash~88 has present value~80, so the required
allowance is $100-80=20$. Debit the old allowance~35 and credit impairment
reversal~35. Cash~100 and net loan~80 now equal deposits~180 and equity~0.
Neither this reversal nor the new payment terms establish recovery from
credit impairment. For a loan that remains recognised, the comparison for a
significant increase in credit risk retains the original, unmodified
origination terms \citep[paragraph~5.5.12 and B5.5.27]{aasb9dated}.

A distinct, ordered credit review at that same close changes the supplied
full-payment probability to $9/10$. Expected~99 has present value~90; allowance
falls~20 to~10, adding an impairment reversal~10 and leaving equity~10.
Compare two histories at this point. In one, reinstatement of the borrower's
operating licence and verified repayment capacity support a finding that the
asset is no longer credit-impaired. The improvement is objectively related to
an event after credit impairment. A significant increase in risk relative to
origination nevertheless remains. In the other history, the supplied credit
assessment still finds impairment despite the same improved cash expectation.
A forecast mean alone does not decide between these findings.

The distinction changes the next year's income presentation. For the cured
asset, paragraph~5.4.2 returns subsequent interest to gross~100: debit gross
loan~10 and credit interest~10. At maturity, the expected shortfall~11 is no
longer discounted, so debit impairment expense~1 and credit allowance~1.
For the still-impaired asset, paragraph~5.4.1(b) applies the original rate to
net~90: debit gross loan~10, credit interest~9 and credit allowance~1.
Both histories reach gross~110, allowance~11 and net~99. Interest is~10 in one
history and~9 in the other, with the compensating difference in impairment
expense \citep[paragraphs~5.4.1--2]{aasb9dated}.

The borrower then pays~110, an outcome admitted by both supplied laws.
Debit cash~110 and allowance~11, credit gross~110 and impairment reversal~11.
Cash~210 equals deposits~180 plus equity~30. Profit since the opening is~55,
equal to~110 received less the opening net loan~55. Equal closing wealth
therefore cannot detect the wrong interest basis. Returning from lifetime
to twelve-month expected credit losses would require a separate reversal of
the significant-risk finding under paragraphs~5.5.3--12.

Ending contractual rights requires another decision. Suppose a different
loan has gross~100 and allowance~20, and an enforceable agreement releases
all remaining rights for an actual payment~85. There is no retained interest
or replacement instrument. Debit cash~85 and allowance~20, credit gross~100
and gain~5: consideration~85 exceeds net carrying amount~80 by~5.
Without the release, the same payment cannot by itself establish full
derecognition \citep[paragraphs~3.2.3 and~3.2.12]{aasb9dated}.

For a transfer, the bank first identifies whether the unit is the whole
asset or an eligible part, then assesses transfer conditions, risks and
rewards, and control. A fixed-price repurchase can retain the risks rather
than establish a sale. The issuer-liability ten-percent comparison does not
answer these asset questions \citep[paragraphs~3.2.2--7]{aasb9dated}.
A replacement asset needs fresh classification, initial measurement and a
purchased-or-originated credit-impairment assessment; a distressed replacement
is not automatically an ordinary new loan with twelve-month losses
\citep[B5.5.25--26]{aasb9dated}. Partial transfers and continuing involvement
use paragraphs~3.2.13--23 \citep[3.2.13--23]{aasb9dated}; the later
revolving-capacity and transfer section develops those conditional methods.

\section{Collateral received after maturity}\label{sec:bank-late-collateral}

Collateral changes expected receipts through its amount, timing and
enforcement costs. Consider a separate loan advancing~100 on~1~January~2024
and requiring~121 on~1~January~2025. Its contractual half-year accumulation
factor is~1.1, equivalent to~21\% effective annual interest. This rate
convention differs from the preceding annual~10\% loan.

The supplied initial law assigns probability $3/5$ to full payment at maturity
and $2/5$ to a default first observed within twelve months. In the default
branch, collateral receipts arrive on~1~July~2025, one half-year after maturity.
Possible net receipts~121, 66.55 and~12.10 have conditional probabilities
$1/4,1/2,1/4$, respectively. Their conditional mean is~66.55. The bank has
supported reasonable prospects of fuller recovery before the sale; 66.55 is
not a known upper limit. These probabilities are illustrative assumptions.

The expected loss compares the present value of contractual cash with that
of all expected receipts. The full-payment value is $121/1.1^2=100$; the
default-branch mean has value $66.55/1.1^3=50$. Thus
\[
 \operatorname{ECL}_0
 =100-\left(\frac{3}{5}\,\frac{121}{1.1^2}
              +\frac{2}{5}\,\frac{66.55}{1.1^3}\right)=20.
\]
The extra half-year belongs in the discounting. Omitting receipts after
maturity produces the wrong loss~40. B5.5.55 expressly includes collateral
cash beyond maturity and deducts obtaining and selling costs; B5.5.43
limits twelve-month losses by the default date, not by the recovery date.
Here all possible defaults occur in the first year, so twelve-month and
lifetime losses coincide \citep[B5.5.43--44 and B5.5.55]{aasb9dated}.

Start with cash~200, deposits~180 and equity~20. Advance~100 and recognise
allowance~20. On~1~July~2024, first accrue gross interest~10. Rolling the
original expectation forward gives net~88 against gross~110, so allowance
rises~20 to~22 through impairment expense~2. A separately observed qualitative
credit-impairment event then establishes the collateral-only expectation:
$66.55/1.1^2=55$. Increase allowance~22 to~55 through impairment expense~33.
The event is an admitted financial-distress finding, not an invented missed
instalment before any instalment is due.

During the next half-year, original-rate interest on net~55 is~5.5.
Debit gross loan~11, credit interest~5.5 and allowance~5.5. At contractual
maturity gross~121 less allowance~60.5 leaves net~60.5. There is no contractual
post-maturity interest, unamortised fee or premium. The initial~100-to~121
difference has been fully amortised. Over the final half-year the expected
receipt value rises~60.5 to~66.55, giving net interest~6.05. Under these terms,
gross remains~121 and allowance falls~60.5 to~54.45: debit allowance~6.05 and
credit interest~6.05. This follows from the Appendix~A definitions and
B5.5.33's gross-minus-discounted-receipts measurement. It is a local derivation
for these terms, not a rule that every overdue contract stops accruing
contractual interest \citep[Appendix~A, paragraph~5.4.1(b), B5.5.33]{aasb9dated}.

The middle recovery occurs: an enforcement agent obtains~79.86, deducts
costs~13.31 already included in the forecast, and remits~66.55.
Debit cash~66.55 and credit gross loan~66.55. Gross and allowance now both
equal~54.45. At this sale, a separate supported exhaustion finding establishes
that no further recovery is reasonably expected. Paragraph~5.4.4 requires
write-off: debit allowance~54.45 and credit gross~54.45. An earlier such
finding would require an earlier write-off; settlement is not a general
waiting rule \citep[paragraph~5.4.4]{aasb9dated}.

Final cash~166.55 equals deposits~180 plus equity $-13.45$. Total interest~21.55
less net impairment expense~55 gives loss~33.45, also equal to~100 advanced
less~66.55 received. The unpaid legal claim~54.45 remains in the contract
history because no release has been supplied. Accounting write-off does not
by itself waive that claim. The bank never recognised the agent's collateral
as a separate owned asset, and never counted its proceeds a second time.
 
